% !PRINTED 381
\weberpage{381}
%<<CONT
$4^{\text{ten}}$ Grades, $F[x, \Theta(x-\xi)]$ ein vollständiges Quadrat sei. Dies giebt zwei Gleichungen zwischen $\xi$ und $\Theta$, aus denen man durch Elimination von $\Theta$ die Doppeltangentengleichung erhält. Da zu jedem $\xi$ nur ein Werth von $\Theta$ gehört (so lange sich nicht zwei Doppeltangenten auf der Linie $L$ schneiden), so kann $\Theta$ rational durch $\xi$ ausgedrückt werden, und zwar in einer Form
\weq{2}{\Theta = \varphi\,(\xi),}
die für jedes zusammengehörige Paar $\xi$, $\Theta$ gilt.

Die Wurzeln der Doppeltangentengleichung ordnen sich ebenso, wie die entsprechenden Doppeltangenten in Complexe, Siebener-Systeme u.\,s.\,w. Wir bezeichnen diese Wurzeln ebenso wie die Doppeltangenten in \S.\,99 durch das Symbol $[ik]$, worin $i$, $k$ die Paare der Ziffern 1 bis 8 durchlaufen und $[ik] = [ki]$ ist.

Betrachten wir irgend zwei von den Doppeltangenten, $p$, $q$, als bekannt, so können wir auf rationalem Wege die Gleichung $u = 0$ eines Kegelschnittes daraus ableiten, der durch die vier Berührungspunkte dieser Doppeltangenten geht, und wir können also die Gleichung \S.\,96, (2)
\[
f = p\,q\,v - u^2
\]
in rationaler Form aufstellen. Dann giebt es nach \S.\,96, (4) unter der Kegelschnittschaar
\[
v + 2\,\lambda\,u + \lambda^2\,p\,q = 0
\]
fünf, die in ein Linienpaar zerfallen, und wenn wir das Product
\[
\prod_{\lambda} (v + 2\,\lambda\,u + \lambda^2\,p\,q) = \Phi
\]
bilden, über die fünf Wurzeln der Gleichung $5^{\text{ten}}$ Grades, von der $\lambda$ abhängt [\S.\,96, (6)], so ist dies Product gleichfalls rational durch die Co\"efficienten von $f$ und von $p$, $q$ ausdrückbar. Dies Product ist aber eine Form $10^{\text{ten}}$ Grades, die in zehn lineare Factoren zerfällt, die mit $p\,q$ zusammen einen \gesperrt{Steiner'schen} Complex bilden.

Die Co\"efficienten in $\Phi$ können nun auch rational durch die beiden den $p$, $q$ entsprechenden Wurzeln $\xi_1$, $\xi_2$ der Doppeltangentengleichung ausgedrückt werden, und wenn wir dann die Abscissen der Schnittpunkte der Linie $L$ mit $\Phi = 0$ aufsuchen, so erhalten wir eine Gleichung $10^{\text{ten}}$ Grades für $\xi$
\[
X(\xi_1, \xi_2, \xi) = 0,
\]
%CONT>>
